\chapter{Research design and methods}
\label{ch:methods}
\section{Design and rationale}
Explain why your design can answer the research question. Identify the unit of
analysis, comparison conditions and controls. Where applicable, describe
recruitment, sampling and inclusion or exclusion criteria.
\section{Materials and procedure}
Document the materials, tools and versions used. Describe each stage in enough
detail for another researcher to reproduce it.
\begin{table}[htbp]
\centering
\caption{Example study protocol; adapt to your research.}
\begin{tabularx}{\linewidth}{@{}lX@{}}
\toprule
\textbf{Stage} & \textbf{Record and report}\\
\midrule
Preparation & Task specification, materials and selection criteria.\\
Collection & Conditions, observations and any protocol deviations.\\
Analysis & Metrics, uncertainty, exclusions and analysis code.\\
\bottomrule
\end{tabularx}
\end{table}
\section{Measures and analysis plan}
Define outcomes, units and the method of comparison before reporting results.
For completion time, one possible summary is
\begin{equation}
\Delta_t = \frac{t_{\mathrm{baseline}}-t_{\mathrm{revised}}}{t_{\mathrm{baseline}}}\times100\%.
\end{equation}
Explain your choice of statistical or qualitative analysis and the treatment of
missing observations. Distinguish confirmatory from exploratory analyses.
\section{Ethics and reproducibility}
Describe consent, privacy, data handling and any required approvals. State how
materials, data and code will be shared, or explain legitimate access restrictions.

